\section*{C7: Programare funcţională}

\subsection{Exercițiul 1 - \( \instr{iszero} \)}
\begin{center}
	\textbf{iszero}:= $\lambda n.\lambda x.\lambda y.((n(\lambda z.y))x)$.
\end{center}
Exemplu:
\begin{align*}
	\textbf{iszero } C_0 & = \textbf{iszero} (\lambda f.\lambda x.x)                               \\
	                     & \to_\beta \lambda x.\lambda y.(((\lambda f.\lambda x.x)(\lambda z.y))x) \\
	                     & \to_\beta \lambda x.\lambda y.((\lambda x.x)x)                          \\
	                     & \to_\beta \lambda x.\lambda y.x= \textbf{true}.                         \\
\end{align*}
similar
\begin{align*}
	\textbf{iszero } C_1 & = \textbf{iszero} (\lambda f.\lambda x.fx)                               \\
	                     & \to_\beta \lambda x.\lambda y.(((\lambda f.\lambda x.fx)(\lambda z.y))x) \\
	                     & \to_\beta \lambda x.\lambda y.((\lambda x.(\lambda z.y)x)x)              \\
	                     & \to_\beta \lambda x.\lambda y.((\lambda z.y)x)                           \\
	                     & \to_\beta \lambda x.\lambda y.y = \textbf{false}.
\end{align*}
scris mai scurt, $\textbf{iszero}\ C_1 \to_\beta^* \textbf{false}$.

\subsection{Succesorul}
dat fiind ca, pentru orice n:

\begin{center}
	$C_n(f)(x)=f^{(n)}(x)$
\end{center}
\begin{center}
	$C_{n+1}(f)(x)=f^{(n+1)}(x)=f(C_n(f)(x)),$
\end{center}

putem defini un $\lambda$ termen care sa reprezinte functia succesor, prin
\begin{center}
	\textbf{succ}:=$\lambda n.\lambda f.\lambda x.(f((nf)x))$.
\end{center}
Exemplu:
\begin{align*}
	\textbf{succ } C_0 & = \textbf{succ} (\lambda f.\lambda x.x)                        \\
	                   & \to_\beta \lambda f.\lambda x.(f(((\lambda f.\lambda x.x)f)x)) \\
	                   & \to_\beta \lambda f.\lambda x.(f((\lambda x.x)x))              \\
	                   & \to_\beta \lambda f.\lambda x.(fx) = C_1
\end{align*}

